\documentclass[french]{book}

\usepackage[utf8]{inputenc}
\usepackage[T1]{fontenc}
\usepackage{lmodern}
\usepackage{geometry}
\usepackage{graphicx}
\usepackage{babel}
\usepackage{amsmath}
\usepackage{amsthm}
\usepackage{amssymb}
\usepackage{hyperref}
\usepackage{stmaryrd}
\usepackage{enumitem}
\usepackage{tcolorbox}
\usepackage{dsfont}
\usepackage{multirow}
\usepackage{etoc}
\usepackage{titlesec}
\usepackage{esint}
\usepackage{cancel}
\usepackage{mathdots}

\setlist[itemize]{label=$\bullet$}


\renewcommand{\P}{\mathbb{P}}
\newcommand{\N}{\mathbb{N}}
\newcommand{\Z}{\mathbb{Z}}
\newcommand{\Q}{\mathbb{Q}}
\newcommand{\R}{\mathbb{R}}
\newcommand{\C}{\mathbb{C}}
\newcommand{\K}{\mathbb{K}}
\renewcommand{\L}{\mathbb{L}}
\newcommand{\U}{\mathbb{U}}
\newcommand{\st}{^*}
\newcommand{\tq}{\middle|}
\newcommand{\inv}{^{-1}}
\newcommand{\E}{\mathbb{E}}
\newcommand{\F}{\mathbb{F}}
\newcommand{\V}{\mathbb{V}}
\renewcommand{\ker}{\text{Ker}}
\newcommand{\im}{\text{Im}}
\newcommand{\card}{\text{Card}}
\newcommand{\Tr}{\text{Tr}}

\theoremstyle{definition}
\newtheorem{ex}{Exercice}[chapter]

\title{Colles BJ 2025-2026}

\begin{document}

\chapter{Semaine 1}

\section{Lafitte}

\begin{ex}
	Soit $n\in\N\st$. On note $z_1,\ldots,z_n$ les racines de $X^n+1$.
	\begin{enumerate}
		\item Soit $k\in\llbracket0,n\rrbracket$. Déterminer la décomposition en éléments simples de la fraction rationnelle $$F_k=\frac{X^k}{X^n+1}$$
		\item Montrer : $$\forall P\in\C[X],\ XP'(X)=\frac{n}{2}P(X)+\frac{2}{n}\sum_{k=1}^{n}\frac{z_kP(z_kX)}{(z_k-1)^2}$$
		\item Soit $P\in\C[X]$. On note $\lVert P\rVert=\underset{\lvert z\rvert=1}{\sup}\lvert P(z)\rvert$. Montrer : $$\lVert P'\rVert\leqslant n\lVert P\rVert$$
	\end{enumerate}
\end{ex}

\begin{ex}
	On considère des matrices carrées $A_1,\ldots,A_p$ telles que $\displaystyle\sum_{k=1}^{p}A_k=I_n$. Montrer l'équivalence : $$\left(\forall k\in\llbracket1,p\rrbracket,\ A_k^2=A_k\right)\iff \left(\forall (i,j)\in\llbracket1,p\rrbracket^2,\ i\ne j\implies A_iA_j=0\right)$$
\end{ex}

\begin{ex}
	Soit $E$ un $\K$-espace vectoriel ainsi que $H_1$ et $H_2$ deux hyperplans de $E$. montrer que $H_1$ et $H_2$ sont isomorphes.
\end{ex}

\section{Nougayrède}

\begin{ex}
	Soit $E$ et $F$ deux $\K$-espaces vectoriels, avec $E$ de dimension finie. Soit $f_1,\ldots,f_p$ des applications linéaires de $E$ dans $F$. On suppose que : $$\text{rg}(f_1+\cdots+f_p)=\sum_{k=1}^{p}\text{rg}(f_k)$$ Montrer que : $$\text{Im}(f_1+\cdots+f_p)=\bigoplus_{k=1}^p\text{Im}(f_k)$$
\end{ex}

\begin{ex}
	Soit $P\in\C[X]$ non constant et unitaire tel que $P(X+1)$ divise $(X+1)P(X)$. Montrer qu'il existe $n\in\N$ tel que $$P(X)=\prod_{k=0}^{n}(X-k)$$
\end{ex}

\begin{ex}
	Soit $E$ un $\K$-espace vectoriel. Soit $(u,v)\in\mathcal{L}(E)^2$ vérifiant $uv=vu$ et $\ker(u)\cap\ker(v)=\{0\}$. Soit $p\geqslant 2$ et $\lambda_1,\ldots,\lambda_ip$ des scalaires deux à deux distincts. \\ Montrer que $\ker(u-\lambda_1v),\ldots,\ker(u-\lambda_pv)$ sont en somme directe.
\end{ex}

\begin{ex}
	Trouver tous les polynômes $P\in\R[X]$ vérifiant : $$\forall k\in\Z,\ \int_{k}^{k+1}P(t)dt=k+1$$
\end{ex}

\begin{ex}
	Soit $E$ un $\K$-espace vectoriel de dimension finie et $f\in\mathcal{L}(E)$. Soit $k$ le plus petit entier naturel tel que $\text{rg}(f^k)=\text{rg}(f^{k+1})$. \\
	Montrer que la dimension de $\ker(f^k)\cap\text{Im}(f)$ est supérieure ou égale à $k-1$.
\end{ex}

\begin{ex}
	Soit $P\in\R[X]$ de degré $2025$ tel que $$\forall k\in\llbracket0,2025\rrbracket,\ P(k)=\frac{k}{k+1}$$
	Calculer $P(2026)$.
\end{ex}

\chapter{Semaine 2}

\section{Nougayrède}

\begin{ex}
	Soit $A$ et $B$ deux éléments de $\mathcal{M}_n(\R)$. On suppose que $A^TB$ est symétrique et que $\ker(A)\cap\ker(B)=\{0\}$.
	\begin{enumerate}
		\item Soit $\lambda$ et $\mu$ deux réels distincts ainsi que $x\in\ker(A-\lambda B)$ et $y\in\ker(A-\mu B)$. \\
		Montrer que $(Bx\mid By)=0$.
		\item Soit $\lambda_1,\ldots,\lambda_p$ des réels deux à deux distincts. Montrer que les sous-espaces $\ker(A-\lambda_1 B),\ldots,\ker(A-\lambda_p B)$ sont en somme directe.
		\item En déduire qu'il existe $\lambda\in\R$ tel que $A-\lambda B\in\mathcal{GL}_n(\R)$.
	\end{enumerate}
\end{ex}

\begin{ex}
	On pose $$J=\begin{pmatrix}
		0&-I_n\\
		I_n&0
	\end{pmatrix}\in\mathcal{M}_{2n}(\R)$$ Soit $M\in\mathcal{M}_{2n}(\R)$ telle que $M^TJM=J$. Montrer que $\det(M)=1$.
\end{ex}

\begin{ex}
	Soit $A,B\in\mathcal{M}_n(\K)$. On suppose que $B^2=B$. Montrer : $$\text{rg}(AB-BA)\leqslant\text{rg}(AB+BA)$$
\end{ex}

\begin{ex}
	Soit $A\in\mathcal{M}_n(\K)$ et $$F=\left\{B\in\mathcal{M}_n(\K)\ : \ AB=BA=0\right\}$$ Déterminer $\dim(F)$ en fonction de $n$ et $r:=\text{rg}(A)$.
\end{ex}

\begin{ex}
	Soit $E$ un $\C$-espace vectoriel de dimension $n$.
	\begin{enumerate}
		\item On suppose qu'il existe $u$ et $v$ deux symétries de $E$ vérifiant $uv+vu=0$. montrer que $n$ est pair et qu'il existe une base $\mathcal{B}$ de $E$ telle que : $$M_\mathcal{B}(u)=\begin{pmatrix}
			I_p&0\\
			0&-I_p
		\end{pmatrix}\ \ \text{et}\ \ M_\mathcal{B}(v)=\begin{pmatrix}
		0&I_p\\
		I_p&0
		\end{pmatrix}$$
		\item Soit $k\in\N$. montrer que les deux assertions suivantes sont équivalentes :
		\begin{enumerate}[label=\roman*]
			\item Il existe $s_1,\ldots,s_k$ des symétries de $E$ vérifiant : $$\forall (i,j)\in\llbracket1,k\rrbracket,\ i\ne j\implies s_is_j+s_js_i=0$$
			\item $k\leqslant 1+2v_2(n)$
		\end{enumerate}
	\end{enumerate}
\end{ex}

\begin{ex}
	Soit $A,B,C,D\in\mathcal{M}_n(\R)$. on suppose que $AC-BD=I_n$ et $BC+AD=0$.\\ Montrer que $\det(AC)\geqslant0$.
\end{ex}

\chapter{Semaine 3}

\section{Nougayrède}

\begin{ex}
	Soit $A\in\mathcal{M}_n(\K)$ et $M=\displaystyle\begin{pmatrix}
		I_n&0\\
		A&A
	\end{pmatrix}$. Montrer que $M$ est diagonalisable si, et seulement si, $A$ est diagonalisable et $A-I_n$ est inversible.
\end{ex}

\begin{ex}
	Soit $E$ un $\K$-espace vectoriel de dimension finie non nulle et $u\in\mathcal{L}(E)$. \\ Déterminer une condition nécessaire et suffisante sur le polynôme minimal de $\pi_u$ pour que pour tout $P\in\K[X]$, l'endomorphisme $P(u)$ soit nul ou inversible.
\end{ex}

\begin{ex}
	Soit $A$ et $B$ dans $\mathcal{M}_n(\R)$ diagonalisables.
	\begin{enumerate}
		\item Montrer que $\R[A]=\R[A^3]$.
		\item On suppose $A^3=B^3$. Montrer $A=B$.
	\end{enumerate}
\end{ex}

\begin{ex}
	Soit $E$ un $\C$-espace vectoriel non nul de dimension finie, ainsi que $u$ et $f$ des endomorphismes de $u$. on se donne $n$ un entier supérieur ou égal à $2$. On suppose que $f$ est diagonalisable à valeurs propres simples, et que $uf-fu=u^n$. \\  Montrer que $u=0$.
\end{ex}

\begin{ex}
	Soit $n\geqslant 2$.
	\begin{enumerate}
		\item Soit $r\in\llbracket1,n\rrbracket$. Expliciter une matrice de rang $r$ et non diagonalisable.
		\item Soit $M\in\mathcal{M}_n(\K)$. On suppose que pour toute matrice $P\in\mathcal{GL}_n(\K)$, $PM$ est diagonalisable. \\ Montrer que $M=0$.
	\end{enumerate}
\end{ex}

\begin{ex}
	Soit $A\in\mathcal{M}_n(\C)$.
	\begin{enumerate}
		\item On suppose que $A$ est diagonalisable. Montrer que la $\C$-algèbre $\C[A]$ est isomorphe à l'algèbre $\C^m$ (pour les opérations termes à termes) pour un certain entier naturel $m$.
		\item Montrer la réciproque.
	\end{enumerate}
\end{ex}

\chapter{Semaine 4}

\section{Lafitte}

\section{Morlot}

\section{Nougayrède}

\begin{ex}
	Soit $E$ un $\C$-espace vectoriel de dimension finie ainsi que $u\in\mathcal{L}(E)$. \\ Montrer que $u$ n'est pas diagonalisable si, et seulement si, il existe $P$ un plan stable par $u$ et une base de $P$ dans laquelle la matrice de l'endomorphisme induit par $u$ sur $P$ est de la forme $\displaystyle\begin{pmatrix}
		\lambda&1\\
		0&\lambda
	\end{pmatrix}$.
\end{ex}

\begin{ex}
	Soit $A\in\mathcal{M}_n(\C)$. On suppose $\Tr(A^n)\ne0$ et $\forall k\in\llbracket1,n-1\rrbracket,\ \Tr(A^k)=0$. \\ Montrer que $A$ est inversible et diagonalisable.
\end{ex}

\begin{ex}
	Soit $A\in\mathcal{M}_n(\C)$. Montrer que $\chi_A$ est scindé à racines simples si, et seulement si, toute matrice nilpotente qui commute avec $A$ est nulle.
\end{ex}

\begin{ex}
	Soit $E$ un $\C$-espace vectoriel de dimension $2$ ainsi que $u$ et $v$ deux automorphismes de $E$. On suppose que $w=uvu\inv v\inv$ commute avec $u$ et $v$. \\ Montrer que $uv=vu$ ou $uv=-vu$.
\end{ex}

\begin{ex}
	Soit $\varphi$ définie sur $\mathcal{M}_2(\C)$ par $\varphi(M)=M^2$. Expliciter l'image de $\varphi$.
\end{ex}

\begin{ex}
	Soit $E$ un $\K$-espace vectoriel de dimension finie et $u\in\mathcal{L}(E)$. Montrer que les deux propriétés suivantes sont équivalentes : 
	\begin{enumerate}[label=\roman*]
		\item $u$ est diagonalisable
		\item $\chi_u$ est scindé et pour tout $P\in\K[X]$, $\ker(P(u))$ et $\text{Im}(P(u))$ sont supplémentaires.
	\end{enumerate}
\end{ex}

\section{Rax}

\section{Vandaële}

\chapter{Semaine n}

\section{Lafitte}

\section{Morlot}

\section{Nougayrède}

\section{Rax}

\section{Vandaële}

\section{Vogel}

\end{document}